\classicsectiondivider{BIG QUESTION: Is the current Vietnamese baseline reliable enough to be deployed as a functional demo?}

\subsection{RQ1: Degradation}

\begin{frame}{RQ1: How Does Missing Video Affect Recognition Accuracy?}
  \small
  \begin{beamercolorbox}[wd=\textwidth,sep=0.70em,center,rounded=true]{block body}
    \textbf{Question:} How does recognition accuracy change as longer or more frequent visual gaps are introduced under laptop-recording conditions?
  \end{beamercolorbox}
  \vspace{0.16cm}

  \begin{columns}[T,onlytextwidth]
    \begin{column}{0.48\textwidth}
      \begin{block}{Evaluation setup}
        \begin{itemize}
          \item Keep the checkpoint, tokenizer, audio, and reference transcripts fixed.
          \item Apply paired deterministic interval masks at increasing visual-dropout ratios.
          \item Measure WER/CER relative to clean AV and repeat across multiple seeds.
        \end{itemize}
      \end{block}
    \end{column}
    \begin{column}{0.48\textwidth}
      \begin{block}{What we want to learn}
        \begin{itemize}
          \item Determine whether degradation is gradual or cliff-like.
          \item Is the released checkpoint already robust enough, or is fine-tuning needed?
        \end{itemize}
      \end{block}
    \end{column}
  \end{columns}
  \note{[Sources] configs and outputs for the paired full-test visual-dropout protocol. No numerical result is shown here; the Implementation section answers this question.}
\end{frame}

\subsection{RQ2: Strategy}

\begin{frame}{RQ2: How Should We Handle Missing Visual Frames?}
  \scriptsize
  Let $v_t$ be the visual input and $m_t\in\{0,1\}$ its availability ($1$: available, $0$: missing) at time step $t$.\\
  Write $\widetilde v_t=m_tv_t$, $f_t^a=E_a(a_t)$, and $f_t^v=E_v(\widetilde v_t)$, using the same encoders as in Background.
  \vspace{0.03cm}

  \renewcommand{\arraystretch}{1.20}
  \begin{tabularx}{\textwidth}{
    >{\bfseries\raggedright\arraybackslash}p{0.18\textwidth}
    >{\centering\arraybackslash}p{0.31\textwidth}
    >{\raggedright\arraybackslash}X}
    \toprule
    Strategy & Fused input $x_t$ & Runtime interpretation \\
    \midrule
    Corrupted AV &
    $[f_t^a;\,f_t^v]$ &
    Zero the missing frames and let the visual encoder process the corrupted sequence. \\
    Interval-gated AV &
    $[f_t^a;\,m_tf_t^v]$ &
    Zero the missing frames and suppress their visual features before fusion. \\
    Audio-only fallback &
    $[f_t^a;\,\mathbf{0}]\ \forall t$ &
    Discard the visual stream for the whole utterance and decode from audio only. \\
    \bottomrule
  \end{tabularx}
  \vspace{0.06cm}

  \begin{columns}[T,onlytextwidth]
    \begin{column}{0.48\textwidth}
      \begin{block}{Evaluation setup}
        Use the same utterance, audio, missing-frame mask, and dropout ratio for every strategy. Compare WER, CER, and inference time.
      \end{block}
    \end{column}
    \begin{column}{0.48\textwidth}
      \begin{alertblock}{What we want to learn}
        When visual input is unreliable, is it better to explicitly suppress visual features at missing time steps, or let the model infer from the corrupted input and the remaining audio-visual context?
      \end{alertblock}
    \end{column}
  \end{columns}
  \note{[Sources] AV-HuBERT.pdf, Section 3.3 and Equations 2-4; src/viavsr/inference/recognition.py. Interval gating is a project-side runtime intervention, not the original AV-HuBERT modality-dropout mechanism. No benchmark result is shown here. Notation follows Background: audio-first concatenation, v for visual input, and a tilde for modified input. Here the modification is runtime zeroing, not the pretraining replacement operation. The expression E_v(tilde v_t) is time-indexed shorthand: the visual frontend processes the sequence and can use neighboring frames, so a zero input frame need not yield a zero feature. The binary availability m_t is distinct from the sequence-level modality-dropout operator D_m.}
\end{frame}
